\documentclass[11pt,a4paper]{article}
\usepackage[margin=2.4cm]{geometry}
\usepackage{amsmath,amssymb,booktabs,array,xcolor}
\usepackage[colorlinks=true,linkcolor=blue!60!black,urlcolor=blue!60!black]{hyperref}
\newcommand{\Lc}{\textsf{\textbf{(c)}}\ }
\newcommand{\Ld}{\textsf{\textbf{(d)}}\ }
\newcommand{\scr}{\texttt{audit\_\allowbreak phase2h\_\allowbreak foundations.py}}
\newcommand{\Li}{\operatorname{Li}}
\newcolumntype{L}[1]{>{\raggedright\arraybackslash}p{#1}}
\emergencystretch=2em
\title{Independent audit of Phase 2h:\\physical phase conventions and the Toda expansion}
\author{Quantum Boussinesq solitons project}
\date{4 November 2026}
\begin{document}
\maketitle
\begin{abstract}
The proposed Toda expansion remains worthwhile, but its numerical targets and
the all-loop master-function interpretation require qualifications. \Ld
Independent phonon quadrature detects an orientation-dependent discrepancy
between the Phase 2h master functional and the Phase 2g physical raw integrals.
The claimed exchange antisymmetry of the raw phase is false in the regular
right-moving sector. The saved edge-extrapolation alternatives do not resolve
even the sign of the quoted subleading Toda coefficients. Conditional on the
stated quasi-momentum shift, the classical Taylor remainder contains
$(\log Y)^2/Y$, whose cancellation in the full phase is unproved. \Lc
We also correct the separated head-on kernel asymptotics and a displayed
soliton profile. Existing calculations are left intact; this note supplies
independent checks and a handoff for concurrent work.
\end{abstract}

\section{Scope, conventions, and reproducibility}
Labels follow the repository convention: (a) established in the literature,
(b) heuristic in the literature, (c) derived or checked in this project,
and (d) conjecture or research judgment. No new external literature claims
are used here. All mathematical checks in this note are implemented in
\texttt{phase2/scripts/}\scr, with results in
\texttt{audit\_phase2h\_results.json}. The script uses 70-digit mpmath arithmetic
and exact SymPy identities. Its JSON records SHA-256 hashes of the source
documents, implementations, and saved data used, and it verifies that these
inputs have not changed during execution. The checkout at the start was
\texttt{0440ed6d345f7a8e94468b16209bc90b92bcdc6d}.

We retain the good-equation convention
\[
 u_{tt}=u_{xx}-u_{xxxx}-3(u^2)_{xx},\qquad
 u=2(\log\tau)_{xx},\qquad k=\sin\alpha,\quad v=\epsilon\cos\alpha.
\]
The phase audit is in the regular region $0<\alpha_j<\pi/3$; its directly
integrated sample has $(\alpha_1,\alpha_2)=(0.3,0.7)$. No new phase convention,
regularisation, or definition of physical charge is selected. The Toda
analysis reuses the saved Phase 2f data and does not claim a new continuum
limit or a fresh high-resolution Bethe-ansatz solve.

\paragraph{What survives.}
The two existing Phase 2h scripts both finish with
\texttt{ALL CHECKS PASSED}. Their algebraic coboundary tests and the
Legendre-transform identity for the proposed collective-coordinate momentum
are reproduced. \Lc\ This audit distinguishes those identities from the
physical identification of the quantities being compared.

\section{The master functional is not yet the physical phase in all orientations}
\subsection{Independent phonon quadrature}
For $\phi_j=\alpha+(2\pi/3,-2\pi/3,0)_j$ and $c_j=(1,1,-2)_j$, we implement
the Phase 2e elementary phase independently:
\begin{align*}
 \theta(q;\alpha,\epsilon)
 &=-\sum_j c_j\left[\operatorname{sgn}(q)\arctan\frac{\sin\phi_j}{|q|}
 +\epsilon\arctan\frac{\sin\phi_j\cos\phi_j}
 {|q|\sqrt{1+q^2}+q^2+\sin^2\phi_j}\right],\\
 I^{\rm raw}_{12}&=\frac1{4\pi}\int_{\mathbb R\setminus\{0\}}
             \theta_1(q)\theta'_2(q)\,dq.
\end{align*}
The phase tends to zero at each infinity, and the integral is split at zero.
These are the existing raw-integral conventions, not a choice of physical
exchange normalization. The independently implemented derivative and
quadrature agree with the Phase 2g closed form to better than $10^{-60}$ at
the four sampled orientations. \Lc\ (\scr, check 1.)

\begin{center}\small
\begin{tabular}{@{}ccrrr@{}}\toprule
$\epsilon_1$ & $\epsilon_2$ & Direct raw integral & Phase 2h master
 & $(\text{master}-\text{raw})/\pi$\\\midrule
$+1$ & $+1$ & $-1.10381243524$ & $-1.10381243524$ & $0$\\
$+1$ & $-1$ & $+0.49133439704$ & $+0.49133439704$ & $0$\\
$-1$ & $+1$ & $-0.49133439704$ & $-3.63292705063$ & $-1$\\
$-1$ & $-1$ & $+1.10381243524$ & $-2.03778021835$ & $-1$\\\bottomrule
\end{tabular}
\end{center}

The same closed-form mismatch is checked at two additional angle pairs.
For a defect with $q_D=0.4$, $\eta=+1$ and $\alpha=0.3$, direct quadrature
also agrees with Phase 2g, whereas the Phase 2h first-difference functional
exceeds that raw soliton--defect integral by $\pi$ for $\epsilon=-1$ and agrees
for $\epsilon=+1$. This example lies in the previously identified stable
defect sector. \Lc

\paragraph{Interpretation.}
The Phase 2h scripts compare a master-function difference to a newly defined
ratio-kernel functional. Those comparisons can pass while the functional
differs from the previously computed physical integral. The identification
across orientations has therefore failed; the numerical algebraic
coboundary identities themselves have not been disproved. \Lc

\subsection{Exchange constants cannot be dropped}
For two right-moving regular solitons, the raw integral obeys
\begin{equation}
 I^{\rm raw}_{12}+I^{\rm raw}_{21}=-\pi,\qquad
 I^{\rm raw}(\alpha,\alpha)=-\frac\pi2.
 \label{eq:raw}
\end{equation}
These statements are independently checked in \scr. \Lc\ They agree with
the integration-by-parts discussion in \texttt{oneloop\_phase.tex}, but
contradict the master report's unqualified claim of antisymmetry and the
paragraph in \texttt{oneloop\_defect\_phase.tex} claiming that the raw phase
vanishes on the diagonal. Adding $\pi/2$ gives an antisymmetric expression
in this right-moving regular sector; it changes the bare triangle identities.
This is a consistency problem to resolve, not a justification for selecting
either normalization here. \Ld

There is also no implication from oddness of a general kernel to exchange
antisymmetry: the already-tested $K(z)=z/(1+z^2)$ yields
$\mathcal I_{SS}(1,2)=\mathcal I_{SS}(2,1)=2.29578918683$ at the same angles.
The independent check reproduces that nonzero symmetric value. \Lc
Consequently, fusion within the algebraic representation does not by itself
select a physical all-loop scattering amplitude. \Ld

\section{What the saved Toda data establish}
\subsection{The quoted subleading coefficients are not resolved}
The coefficient used in the proposed next step is obtained from only two
edge-extrapolated values:
\begin{equation}
 \widehat c_1(r)=32\,[R_0(16,r)-R_0(32,r)].
 \label{eq:c1}
\end{equation}
This is the slope under an assumed linear dependence on $1/Y$, not an
independent extraction of an asymptotic coefficient. \Lc\ The saved
\texttt{matching\_analysis.json} already contains alternative edge fits.
Applying \eqref{eq:c1} to each yields the following checked results
(\scr, check 4):
\begin{center}\small
\begin{tabular}{@{}crrrr@{}}\toprule
$r$ & Quadratic edge fit & Cubic edge fit & $h\log h$ edge fit
 & Propagated spread\\\midrule
0.4 & $-0.007080$ & $-0.007261$ & $+0.009148$ & $0.018275$\\
0.5 & $-0.020991$ & $-0.021363$ & $+0.017819$ & $0.045090$\\
0.6 & $-0.048219$ & $-0.048907$ & $+0.030317$ & $0.093195$\\\bottomrule
\end{tabular}
\end{center}
Here the last column is $32(s_{16}+s_{32})$, with $s_Y$ the recorded edge-fit
spread. It is a sensitivity diagnostic, not a rigorous or statistical error
bar. The sign change in the preceding columns uses a common alternative fit
at both $Y$ values, so it does not depend on treating their errors as
independent. \Lc\ Better asymptotic control could eventually exclude that
alternative, but the present data do not do so. \Ld

\subsection{The proposed leading limit remains a candidate}
The central large-$Y$ intercepts lie within $1.3\times10^{-5}$ to
$7.1\times10^{-5}$ of the proposed
\[
 c_0^*(r)=-\frac4\pi\big(\chi_2(r)-r\big).
\]
This numerical proximity is reproduced. \Lc\ However, the propagated
edge-fit spreads are $0.00150$, $0.00369$, and $0.00761$, respectively,
more than fifty times those discrepancies. Agreement of central estimates
does not establish accuracy at the displayed discrepancy scale. \Ld

In \texttt{charge\_and\_frame\_audit.py}, the statement
\begin{quote}\small\ttfamily
dPhi\_loop\_asymp = R0\_target - dPhi\_cl\_asymp
\end{quote}
builds the proposed answer into the loop asymptotic. The subsequent check
is a numerical compatibility check, not an independent analytical
derivation of the answer. \Lc\ Claim 3.2 of
\texttt{charge\_frame\_master.tex} asserts the corresponding loop expansion
but does not derive it from the Bethe integral equation. This audit neither
proves nor disproves $c_0^*(r)$; an independent edge-equation derivation is
still needed. \Ld

\section{A logarithmic term that the next Toda calculation must track}
Retain the Phase 2h classical edge function and its proposed shifts:
\begin{align*}
 \Phi_{\rm cl}(u_1,u_2)
 &=2(u_1-u_2)\log\frac{1-\sqrt{u_2/u_1}}{1+\sqrt{u_2/u_1}}
       +4\sqrt{u_1u_2},\\
 (y_1,y_2)&=(Y,r^2Y),\qquad
 \delta u_1=-\frac{\log Y}{2\pi}+c,\qquad
 \delta u_2=\delta u_1-\frac{\log r}{\pi}.
\end{align*}
The shift formulas are inputs to this conditional calculation, not newly
derived Bethe-ansatz results. Exact differentiation gives the quadratic
Taylor term \Lc\ (\scr, check 3):
\begin{equation}
 Q=\frac12\delta u^T D^2\Phi_{\rm cl}\,\delta u
   =\frac{(\delta u_2-r^2\delta u_1)^2}{r(1-r^2)Y}.
 \label{eq:Q}
\end{equation}
In particular,
\begin{equation}
 Q=\frac{1-r^2}{4\pi^2r}\frac{(\log Y)^2}{Y}
      +O\!\left(\frac{\log Y}{Y}\right),\qquad 0<r<1.
 \label{eq:logQ}
\end{equation}
SymPy verifies the rank-one Hessian and this coefficient; an independent
70-digit Taylor comparison at $Y=10^{12}$ verifies the quadratic term with
relative error below $10^{-10}$. \Lc

Thus the stated $O(Y^{-1})$ remainder after the first-order classical Taylor
term is not a correct strict asymptotic bound under the proposed logarithmic
shift. \Lc\ This does \emph{not} prove a surviving logarithm in the full
residual: the intrinsic loop phase and further momentum corrections may
cancel it. The next calculation must establish those cancellations before
assuming a pure $c_1(r)/Y$ expansion. \Ld

\section{Additional checks: one correction and one positive control}
\subsection{The separated head-on kernels}
Let $\alpha_1=r\lambda$, $\alpha_2=\lambda$, with directions $(+1,-1)$.
Claim 4.3 of Phase 2h assigns a leading logarithm to each separated kernel.
An independent expansion instead gives \Lc\ (\scr, check 2)
\begin{align}
 I_{\chi_2}^{\rm head}
 &=\frac{9r}{\pi}\lambda^2+O(\lambda^4),\label{eq:chihead}\\
 I_{\ln}^{\rm head}
 &=\frac{r\lambda^2}{\pi}
 \left[-3\log(r\lambda^2)+6\log2-\frac32\log3-12\right]
       +O(\lambda^4|\log\lambda|).
 \label{eq:loghead}
\end{align}
Their sum reproduces the existing total head-on formula, including its
nonzero logarithmic term. The old script checked only that sum; its passing
status did not verify the separately stated asymptotics. \Lc

For a short derivation of \eqref{eq:chihead}, at zero amplitudes the three
half-angle tangents and their derivatives are
$w=(\sqrt3,-\sqrt3,0)$ and $w'=(2,2,1/2)$. For
$K_{\chi_2}=-\operatorname{Re}\chi_2/\pi$, the mixed derivative of
$K_{\chi_2}(w_iw_j)$ is
$-w'_iw'_j/[\pi(1-w_i^2w_j^2)]$. The weighted sum is $9/\pi$, verified
exactly by SymPy. The sum is odd in each amplitude and analytic in a
neighbourhood of zero, giving the displayed remainder. Independent numerical
checks use $r=0.25,0.5,0.75$ and $\lambda=10^{-4},10^{-7},10^{-10}$. \Lc

\subsection{Charges and the profile}
The identity $\Delta P_{\rm can}=dL_1/dv$, using the proposed
$L_1=v\Delta P_{\rm mode}-\Delta E_1^{\rm ren}$, agrees with the existing
closed form to better than $10^{-60}$ at the regular sampled points. \Lc
This is a positive algebraic control. It does not independently derive the
renormalised effective action, or settle which full conserved charges enter
soliton absorption and the Toda frame dictionary. Those physical
identifications remain outside this audit. \Ld

The displayed profile in Phase 2h also has a sign and amplitude error. From
its stated $\tau=1+e^{k(x-vt)}$ convention one obtains
\[
 u=\frac{k^2}{2}\operatorname{sech}^2\frac{k(x-vt)}2,
\]
not $-k^2\operatorname{sech}^2(k(x-vt)/2)$. SymPy verifies the tau derivative.
Substituting the quoted profile in the travelling-wave equation
$u''-k^2u+3u^2=0$ gives residual $9k^4/2$ at its centre. \Lc
This displayed error alone does not invalidate the implemented charge
identities, which are checked separately.

\section{Handoff, claims summary, and failures recorded}
\paragraph{Recommended division of work.}\Ld
The other investigation can proceed with the direct Toda edge expansion.
It should treat $c_0^*$ as a candidate, retain logarithmic subleading terms,
and avoid fitting to the quoted $c_1$ numbers. An all-loop physical kernel
or ring quantisation should explicitly track orientation and zero-mode
constants; the algebraic master functional cannot yet be substituted for
the physical phase without that work. These recommendations select no new
regularisation or normalization on the PI's behalf.

\begin{center}\small
\begin{tabular}{@{}L{6.1cm}L{3.6cm}L{4.5cm}@{}}\toprule
Claim or test & Audit status & Evidence\\\midrule
Physical raw phase equals the master functional in all directions
 & Contradicted \Lc & Check 1: independent quadrature; $\pi$ discrepancies\\
Raw phase is antisymmetric and vanishes on the diagonal
 & Contradicted \Lc & Check 1; Eq.~\eqref{eq:raw}\\
Algebraic coboundary tests
 & Existing tests pass \Lc & Original master script rerun\\
Quoted $c_1$ values are established two-loop targets
 & Unsupported \Ld & Check 4: saved fit alternatives reverse their signs \Lc\\
Proposed exact $c_0^*$ limit
 & Open \Ld & Numerical compatibility; independent derivation missing\\
First-order classical remainder is $O(1/Y)$
 & Contradicted conditionally \Lc & Check 3; Eqs.~\eqref{eq:Q}--\eqref{eq:logQ}\\
Separated head-on logarithms
 & Corrected \Lc & Check 2; total asymptotic survives\\
Legendre-transform momentum formula
 & Confirmed conditionally \Lc & Check 5; $<10^{-60}$ residual\\
Displayed negative soliton profile
 & Contradicted \Lc & Check 5; tau derivative and ODE residual\\\bottomrule
\end{tabular}
\end{center}

\paragraph{Failures recorded.}
\begin{enumerate}
\item The physical-identification, antisymmetry, diagonal-value,
separated-head-on, strict-remainder, and displayed-profile claims fail the
checks detailed above. These failures are retained explicitly as audit
results; no tolerances were relaxed to make those claims pass.
\item The first audit run stalled while simplifying a Hessian written with
nested square roots and was interrupted. Replacing $u_1=s^2$, $u_2=t^2$
and applying the exact chain rule makes the same symbolic Hessian check
tractable; no mathematical assertion is weakened.
\item The tau-profile equality check initially rewrote only one side into
exponentials, leaving a mixed exponential/hyperbolic expression unevaluated
and failing the symbolic-zero assertion. Rewriting both sides consistently
and applying exact rational cancellation proves the same identity. No
numerical tolerance is involved.
\end{enumerate}

\paragraph{Verification scope.}
The standalone audit and both original Phase 2h scripts were run. The
full repository suite was not rerun: the audit changes no existing
implementation, and does not claim to revalidate all earlier phases.
\texttt{ALL CHECKS PASSED} in the new script means that its positive controls
and reproductions of the recorded prior-claim failures passed. It does
not mean the prior claims are now correct. The LaTeX note is built twice
with \texttt{pdflatex} and its rendered pages are inspected.
\end{document}
